\section{Flow-Analyse}
\label{sec:flow}

\subsection{Die Rechnung des Eigent\"umers}

Der Anleger kauft keinen Gewinn, sondern einen Anspruch auf einen Zahlungsstrom. Die
Rechnung dieses Abschnitts stellt deshalb den Kaufpreis (den B\"orsenwert von
\Mio{\Boersenwert}{EUR}) den k\"unftigen Zufl\"ussen gegen\"uber und schlie{\ss}t mit dem
Endwert, dem Buchwert des Eigenkapitals von \Mio{\Eigenkapital}{EUR}.

\annahme{Der freie Zufluss bleibt \"uber den gesamten Horizont konstant, und der Endwert
ist der heutige Buchwert des Eigenkapitals, also ohne Wachstum und ohne Aufschlag. Beides
ist bewusst streng. Ohne Endwert w\"urde die Rechnung zu einer reinen
Aussch\"uttungsrechnung und der Schwellenkurs erheblich zu niedrig ausfallen; mit einem
Endwert \"uber dem Buchwert w\"urde der Bericht genau die Erwartung setzen, die er
pr\"ufen soll.}

\annahme{Drei Szenarien, alle aus der eigenen Reihe 2019--2025 und nicht aus einer
Sch\"atzung: das 25.~Perzentil (\Mio{\FcfPes}{EUR}), der Median (\Mio{\FcfMedian}{EUR})
und das 75.~Perzentil (\Mio{\FcfOpt}{EUR}). Der Zufluss des Gesch\"aftsjahres 2025
(\Mio{\FcfVoll}{EUR}) f\"allt mit dem Median zusammen; das Basisszenario ist damit
zugleich das letzte Ist-Jahr. Die Spanne der Reihe reicht von \Mio{\FcfMin}{EUR}
(\FcfMinJahr) bis \Mio{\FcfMax}{EUR} (\FcfMaxJahr).}

\subsection{Woher der Zufluss kommt und wohin er geht}

Es gibt keine Finanzierungsfrage: GEA weist eine Nettoliquidit\"at von
\Mio{\Nettoliquiditaet}{EUR} aus, kein Zins bindet den Zufluss vorab. Die Frage ist die
Verwendung. 2025 gingen \Mio{\DividendeSumme}{EUR} an die Aktion\"are und
\Mio{\RueckkaufSumme}{EUR} in R\"uckk\"aufe, zusammen \AusschuettungAnteil{} des freien
Zuflusses. Der Rest verblieb im Unternehmen und erh\"ohte die Nettoliquidit\"at.

\bk{Ein Unternehmen, das \AusschuettungAnteil{} seines freien Zuflusses aussch\"uttet und
trotzdem Liquidit\"at aufbaut, verdient mehr, als es f\"ur sein Wachstum braucht. Das ist
die angenehme Lage; sie hat aber eine Kehrseite f\"ur den K\"aufer von heute. Wenn das
Gesch\"aft kein zus\"atzliches Kapital aufnehmen kann, ohne die Rendite zu verw\"assern,
dann ist der Wert der Aktie weitgehend der Barwert der Aussch\"uttungen und nicht der
Barwert eines wachsenden Kapitalstocks. Genau deshalb schl\"agt die Rechnung so hart
gegen den heutigen Kurs aus.}

\subsection{Was der Kaufpreis verzinst}

Der Schwellenkurs ist der Kurs, zu dem der interne Zinsfu{\ss} genau die H\"urde erreicht.

\schwellentabelle

Im Basisszenario betr\"agt er auf zehn Jahre \je{\SchwelleMEDZehn}{EUR}, bei der
vorsichtigeren H\"urde von \Pz{\HuerdeLang} \je{\SchwelleMEDZehnLang}{EUR}. Der heutige
Kurs von \je{\Kurs}{EUR} ist das \VielfachesMED-fache des Schwellenkurses im Basisfall,
im optimistischen Szenario das \VielfachesOPT-fache und im pessimistischen das
\VielfachesPES-fache.

\bk{Dass der Schwellenkurs mit dem Horizont steigt und nicht f\"allt, ist kein Fehler: Je
l\"anger der Horizont, desto mehr Jahreszufl\"usse tr\"agt die Rechnung und desto weniger
wiegt der abgezinste Endwert. Bemerkenswert ist das Verh\"altnis der beiden Spannen: Die
drei Szenarien reichen auf zehn Jahre von \je{\SchwellePESZehn}{EUR} bis
\je{\SchwelleOPTZehn}{EUR}, die Wahl der H\"urde verschiebt den Basisfall von
\je{\SchwelleMEDZehn}{EUR} auf \je{\SchwelleMEDZehnLang}{EUR}. Beide Effekte sind von
derselben Gr\"o{\ss}enordnung, und keiner der sechs Werte kommt dem heutigen Kurs nahe.
Die Frage ist also nicht, welches Jahr der Reihe man f\"ur typisch h\"alt.}
